\subsection{Correlazione tra Valutazione Euristica e Test Utenti}

In questa fase dell'analisi, abbiamo incrociato i dati qualitativi emersi durante le sessioni di test con gli utenti con le criticità strutturali già individuate attraverso la valutazione euristica (Sezione \ref{subsec:valutazione_euristica}). L'obiettivo di questo confronto è verificare in che misura le violazioni dei principi di Nielsen si riflettano sull'esperienza d'uso reale, mettendo in relazione le criticità individuate con le difficoltà, la frustrazione e il disorientamento emersi durante i test con gli utenti.

Un primo elemento di convergenza riguarda l'\textbf{incoerenza linguistica e terminologica}. Durante i test, diversi partecipanti hanno manifestato confusione nell'interpretazione di pulsanti ed etichette, segnalando in particolare l'ambiguità della voce \textit{StudyFlow AI} e la scarsa distinzione tra le azioni \textit{Salva} e \textit{Scarica}. Tali osservazioni confermano quanto emerso dalla valutazione euristica in relazione al Principio 2 di Nielsen (\textit{Adeguare il sistema al mondo reale}): l'utilizzo di una terminologia non coerente, alternando italiano e inglese e adottando nomi poco rappresentativi delle funzionalità, rende più difficile comprendere la struttura dell'applicazione e individuare il percorso corretto per svolgere i compiti.

Per quanto riguarda la \textbf{gestione degli errori e del feedback visivo}, molti utenti hanno riportato un senso di smarrimento quando un'operazione non andava a buon fine, soprattutto in assenza di messaggi che spiegassero chiaramente la causa del problema. A ciò si aggiunge il fastidio generato dai continui \textit{refresh} della pagina, che interrompevano il flusso di interazione e portavano frequentemente a dimenticare informazioni già inserite, come la selezione del corso durante il caricamento di un documento. Questi comportamenti osservati durante i test confermano le criticità individuate rispetto al Principio 1 (\textit{Far vedere lo stato del sistema}), evidenziando come un feedback insufficiente favorisca errori ripetuti e aumenti il carico cognitivo dell'utente.

Il punto di convergenza più evidente e numericamente rilevante riguarda la \textbf{mancanza di consistenza e le difficoltà di navigazione}. La valutazione euristica (Figura \ref{fig:violazioni-nielsen}) aveva infatti evidenziato come criticità assoluta dell'interfaccia originale proprio il \textbf{Principio 4 di Nielsen (Consistenza)}, registrando un picco di ben \textbf{29 violazioni}.

I test sperimentali hanno fornito una puntuale validazione empirica di questo dato teorico, dimostrando come la maggior parte delle problematiche sia riconducibile all'area della \textbf{cognizione} (Figura \ref{fig:categorie-principi-nielsen}), ovvero a problemi che richiedono all'utente uno sforzo maggiore per decodificare il sistema. I frequenti fenomeni di disorientamento osservati nel \textit{Compito C3} (Pianificare un esame) e nello \textit{Scenario S1} derivano direttamente da questa carenza strutturale. Label non coerenti e il posizionamento non standard degli elementi di comando hanno costretto gli utenti a percorsi di navigazione poco intuitivi e a esplorazioni casuali (\emph{trial and error}) per completare operazioni semplici, come collegare un appunto della \textit{community} a un esame.

Nel complesso, i risultati dei test di usabilità confermano le evidenze emerse dalla valutazione euristica. Le violazioni dei principi di Nielsen individuate durante l'ispezione non rappresentano soltanto criticità teoriche, ma trovano riscontro nel comportamento degli utenti, traducendosi in errori, rallentamenti e momenti di incertezza durante l'interazione. Il confronto tra i due metodi rafforza quindi l'affidabilità dei risultati ottenuti e mostra come le problematiche individuate abbiano un impatto concreto sull'usabilità complessiva del sistema.
